% Options for packages loaded elsewhere
\PassOptionsToPackage{unicode}{hyperref}
\PassOptionsToPackage{hyphens}{url}
\PassOptionsToPackage{dvipsnames,svgnames,x11names}{xcolor}
%
\documentclass[
  super,
  preprint,
  5p,
  twocolumn]{elsarticle}

\usepackage{amsmath,amssymb}
\usepackage{iftex}
\ifPDFTeX
  \usepackage[T1]{fontenc}
  \usepackage[utf8]{inputenc}
  \usepackage{textcomp} % provide euro and other symbols
\else % if luatex or xetex
  \usepackage{unicode-math}
  \defaultfontfeatures{Scale=MatchLowercase}
  \defaultfontfeatures[\rmfamily]{Ligatures=TeX,Scale=1}
\fi
\usepackage{lmodern}
\ifPDFTeX\else  
    % xetex/luatex font selection
\fi
% Use upquote if available, for straight quotes in verbatim environments
\IfFileExists{upquote.sty}{\usepackage{upquote}}{}
\IfFileExists{microtype.sty}{% use microtype if available
  \usepackage[]{microtype}
  \UseMicrotypeSet[protrusion]{basicmath} % disable protrusion for tt fonts
}{}
\makeatletter
\@ifundefined{KOMAClassName}{% if non-KOMA class
  \IfFileExists{parskip.sty}{%
    \usepackage{parskip}
  }{% else
    \setlength{\parindent}{0pt}
    \setlength{\parskip}{6pt plus 2pt minus 1pt}}
}{% if KOMA class
  \KOMAoptions{parskip=half}}
\makeatother
\usepackage{xcolor}
\setlength{\emergencystretch}{3em} % prevent overfull lines
\setcounter{secnumdepth}{5}
% Make \paragraph and \subparagraph free-standing
\ifx\paragraph\undefined\else
  \let\oldparagraph\paragraph
  \renewcommand{\paragraph}[1]{\oldparagraph{#1}\mbox{}}
\fi
\ifx\subparagraph\undefined\else
  \let\oldsubparagraph\subparagraph
  \renewcommand{\subparagraph}[1]{\oldsubparagraph{#1}\mbox{}}
\fi


\providecommand{\tightlist}{%
  \setlength{\itemsep}{0pt}\setlength{\parskip}{0pt}}\usepackage{longtable,booktabs,array}
\usepackage{calc} % for calculating minipage widths
% Correct order of tables after \paragraph or \subparagraph
\usepackage{etoolbox}
\makeatletter
\patchcmd\longtable{\par}{\if@noskipsec\mbox{}\fi\par}{}{}
\makeatother
% Allow footnotes in longtable head/foot
\IfFileExists{footnotehyper.sty}{\usepackage{footnotehyper}}{\usepackage{footnote}}
\makesavenoteenv{longtable}
\usepackage{graphicx}
\makeatletter
\def\maxwidth{\ifdim\Gin@nat@width>\linewidth\linewidth\else\Gin@nat@width\fi}
\def\maxheight{\ifdim\Gin@nat@height>\textheight\textheight\else\Gin@nat@height\fi}
\makeatother
% Scale images if necessary, so that they will not overflow the page
% margins by default, and it is still possible to overwrite the defaults
% using explicit options in \includegraphics[width, height, ...]{}
\setkeys{Gin}{width=\maxwidth,height=\maxheight,keepaspectratio}
% Set default figure placement to htbp
\makeatletter
\def\fps@figure{htbp}
\makeatother

\makeatletter
\makeatother
\makeatletter
\makeatother
\makeatletter
\@ifpackageloaded{caption}{}{\usepackage{caption}}
\AtBeginDocument{%
\ifdefined\contentsname
  \renewcommand*\contentsname{Table of contents}
\else
  \newcommand\contentsname{Table of contents}
\fi
\ifdefined\listfigurename
  \renewcommand*\listfigurename{List of Figures}
\else
  \newcommand\listfigurename{List of Figures}
\fi
\ifdefined\listtablename
  \renewcommand*\listtablename{List of Tables}
\else
  \newcommand\listtablename{List of Tables}
\fi
\ifdefined\figurename
  \renewcommand*\figurename{Figure}
\else
  \newcommand\figurename{Figure}
\fi
\ifdefined\tablename
  \renewcommand*\tablename{Table}
\else
  \newcommand\tablename{Table}
\fi
}
\@ifpackageloaded{float}{}{\usepackage{float}}
\floatstyle{ruled}
\@ifundefined{c@chapter}{\newfloat{codelisting}{h}{lop}}{\newfloat{codelisting}{h}{lop}[chapter]}
\floatname{codelisting}{Listing}
\newcommand*\listoflistings{\listof{codelisting}{List of Listings}}
\makeatother
\makeatletter
\@ifpackageloaded{caption}{}{\usepackage{caption}}
\@ifpackageloaded{subcaption}{}{\usepackage{subcaption}}
\makeatother
\makeatletter
\makeatother
\usepackage{float}
\makeatletter
\let\oldlt\longtable
\let\endoldlt\endlongtable
\def\longtable{\@ifnextchar[\longtable@i \longtable@ii}
\def\longtable@i[#1]{\begin{figure}[H]
\onecolumn
\begin{minipage}{0.5\textwidth}
\oldlt[#1]
}
\def\longtable@ii{\begin{figure}[H]
\onecolumn
\begin{minipage}{0.5\textwidth}
\oldlt
}
\def\endlongtable{\endoldlt
\end{minipage}
\twocolumn
\end{figure}}
\makeatother
\journal{Computational Toxicology}
\ifLuaTeX
  \usepackage{selnolig}  % disable illegal ligatures
\fi
\usepackage[]{natbib}
\bibliographystyle{elsarticle-num}
\IfFileExists{bookmark.sty}{\usepackage{bookmark}}{\usepackage{hyperref}}
\IfFileExists{xurl.sty}{\usepackage{xurl}}{} % add URL line breaks if available
\urlstyle{same} % disable monospaced font for URLs
\hypersetup{
  pdftitle={ToxPipe: Leveraging Expert Entrained AI and Cloud Computing for Advanced Computational Toxicology},
  pdfauthor={Trey O. Saddler; Charles Schmitt; David Reif; Scott Auerbach},
  pdfkeywords={retrieval-augmented generation, generative ai, large
language models, toxicology, autonomous agents},
  colorlinks=true,
  linkcolor={blue},
  filecolor={Maroon},
  citecolor={Blue},
  urlcolor={Blue},
  pdfcreator={LaTeX via pandoc}}

\setlength{\parindent}{6pt}
\begin{document}

\begin{frontmatter}
\title{ToxPipe: Leveraging Expert Entrained AI and Cloud Computing for
Advanced Computational Toxicology}
\author[1]{Trey O. Saddler%
\corref{cor1}%
}
 \ead{trey.saddler@nih.gov} 
\author[2]{Charles Schmitt%
%
}
 \ead{charles.schmitt@nih.gov} 
\author[]{David Reif%
%
}
 \ead{david.reif@nih.gov} 
\author[1]{Scott Auerbach%
%
}
 \ead{auerbachs@niehs.nih.gov} 

\affiliation[1]{organization={Divisipn of Translational Toxicology,
National Institutes of Environmental Health Sciences, Predictive
Toxicology Branch},addressline={530 Davis
Dr},city={Durham},postcode={27709},postcodesep={}}
\affiliation[2]{organization={National Institutes of Environmental
Health Sciences, Office of Data Science},addressline={530 Davis
Dr},city={Durham},postcode={27709},postcodesep={}}

\cortext[cor1]{Corresponding author}




        
\begin{abstract}
The Division of Translational Toxicology at NIEHS introduces ToxPipe, a
pioneering project aiming to harness expert entrained AI-based systems
and cloud computing to revolutionize the analysis and interpretation of
toxicological properties of compounds. ToxPipe repurposes advanced AI
models like JARVIS and Auto-GPT to facilitate the exploration of
toxicologically-relevant data through natural language instructions,
enabling a semi-autonomous approach to toxicological characterization
and data integration. This project is grounded in the success of expert
entrained AI models like Auto-GPT and JARVIS, which employ OpenAI's
GPT-based models as controllers to connect fine-tuned, expert AI models,
allowing the execution of complex tasks through plain language
interfaces. ToxPipe utilizes Azure Cloud to enhance the speed and
efficiency of toxicological predictions, leveraging its pay-as-you-go
model for cost-effective utilization of resources and direct access to
crucial OpenAI models. The project also integrates Posit Connect for the
deployment of Shiny web applications, Jupyter notebooks, and APIs,
allowing seamless interaction with ToxPipe through a primary Shiny
application and facilitating the development of additional applications
using the ToxPipe API. The project leverages existing resources on
BioWulf and NIEHS HPC clusters for heavy compute loads, exploring hybrid
deployment models and reducing cloud costs. The utilization of these
resources also allows easy integration of scientific tools and data
already loaded on BioWulf. The implementation of ToxPipe is poised to
significantly accelerate the research process in toxicology, serving as
a blueprint for harnessing cloud computing in scientific discovery. The
project's success will be evaluated against a set of clearly defined
metrics and performance indicators, including accuracy, speed,
cost-effectiveness, and the quality and scope of ToxPipe generated
products, aligning with NIH's strategic goals in data science and
research. In conclusion, ToxPipe represents a groundbreaking initiative
in computational toxicology, combining expert entrained AI, cloud
computing, and hybrid deployments to advance toxicological workflows and
scientific knowledge discovery. The well-defined milestones and
experienced team are set to ensure the project's success and contribute
to the evolution of toxicological research methodologies.
\end{abstract}





\begin{keyword}
    retrieval-augmented generation \sep generative ai \sep large
language models \sep toxicology \sep 
    autonomous agents
\end{keyword}
\end{frontmatter}
    \hypertarget{introduction}{%
\section{Introduction}\label{introduction}}

\begin{itemize}
\tightlist
\item
  Background on toxicological information retrieval.
\item
  The significance of AI in enhancing toxicological data interpretation.
\item
  Objective of the ToxPipe project.
\end{itemize}

Toxicological research is an important field that plays a critical role
in protecting public health and the environment. However, the sheer
volume of toxicological data available can make it difficult for
researchers to quickly and accurately retrieve and interpret the
information they need. ToxPipe is designed to address this challenge by
providing an AI-assisted platform that can help researchers quickly and
accurately retrieve and interpret toxicological information.

\hypertarget{materials-and-methods}{%
\section{Materials and methods}\label{materials-and-methods}}

ToxPipe uses a combination of natural language processing, machine
learning, and data visualization techniques to help researchers retrieve
and interpret toxicological information. The platform is designed to be
user-friendly and intuitive, with a simple and intuitive interface that
allows researchers to quickly and easily search for and retrieve the
information they need.

\hypertarget{data-collection}{%
\subsection{Data Collection}\label{data-collection}}

Description of the data sources and collection methods.

\hypertarget{ai-model-development}{%
\subsection{AI Model Development}\label{ai-model-development}}

\begin{itemize}
\tightlist
\item
  Overview of the AI models used.
\item
  Training and validation of the models.
\end{itemize}

\hypertarget{platform-design}{%
\subsection{Platform Design}\label{platform-design}}

Description of the ToxPipe platform's user interface, features, and
functionalities.

\hypertarget{results}{%
\section{\texorpdfstring{Results }{Results }}\label{results}}

ToxPipe has been tested on a variety of toxicological datasets, and has
been shown to be highly effective at retrieving and interpreting
toxicological information. The platform has also been well-received by
toxicologists and researchers, who have praised its ease of use and
effectiveness.

\hypertarget{model-performance}{%
\subsection{Model Performance}\label{model-performance}}

\begin{itemize}
\tightlist
\item
  Metrics showcasing the performance of the AI models.
\item
  Comparative analysis with existing systems, if any.
\end{itemize}

\hypertarget{platform-usability}{%
\subsection{Platform Usability}\label{platform-usability}}

Feedback and results from user testing of the ToxPipe platform.

\hypertarget{discussion}{%
\section{Discussion}\label{discussion}}

\begin{itemize}
\tightlist
\item
  Interpretation of the results.
\item
  Implications of the findings.
\item
  Potential applications and impact of the ToxPipe platform in the field
  of toxicology.
\end{itemize}

ToxPipe represents a significant advance in the field of toxicological
research, providing researchers with a powerful tool for quickly and
accurately retrieving and interpreting toxicological information. The
platform has the potential to significantly improve the efficiency and
effectiveness of toxicological research, and could have important
implications for public health and the environment.

\hypertarget{conclusions}{%
\section{Conclusions}\label{conclusions}}

Summary of the main findings and future directions for the ToxPipe
project.

\hypertarget{acknowledgements}{%
\section{Acknowledgements}\label{acknowledgements}}

Acknowledging individuals or institutions that contributed to the
project but are not listed as authors.

\hypertarget{declaration-of-generative-ai-and-ai-assisted-technologies-in-the-writing-process}{%
\section{Declaration of Generative AI and AI-assisted technologies in
the writing
process}\label{declaration-of-generative-ai-and-ai-assisted-technologies-in-the-writing-process}}

During the preparation of this work the author(s) used ChatGPT (GPT-4)
in order to generate an initial abstract draft from the contents of a
funding proposal for this project.

Additionally, ChatGPT (GPT-4) was used to generate an initial outline
for this manuscript, which was then heavily edited by the author(s) and
expanded upon.

After using this tool/service, the author(s) reviewed and edited the
content as needed and take(s) full responsibility for the content of the
publication.


  \bibliography{bibliography.bib}


\end{document}
